\chapter{Limitations}
\label{ch:limitations}

This chapter states the limitations of \projname{} honestly. Each
item is either an intentional trade-off, a scope decision, or an
open technical issue. Every claim in this chapter is verifiable from
the current repository.

\section{Data Modality Limitations}
\label{sec:lim-data}

\begin{itemize}[leftmargin=1.4em]
  \item \textbf{Tabular only.} The pipeline assumes rectangular
        pandas-friendly data. Text, images, audio, time series,
        and graphs are out of scope.
  \item \textbf{Single-file datasets.} Multi-table joins, relational
        databases, and streaming data are not supported.
  \item \textbf{Memory-bound.} Datasets are loaded whole with
        \code{pandas.read\_csv} or \code{read\_parquet}. Extremely
        wide or extremely long datasets can exhaust RAM before the
        upload cap is reached.
\end{itemize}

\section{Compute Limitations}
\label{sec:lim-compute}

\begin{itemize}[leftmargin=1.4em]
  \item \textbf{CPU only.} The recommended install path does not
        include GPU-enabled builds of XGBoost or LightGBM. Users can
        install them separately, but no CUDA integration is provided
        out of the box.
  \item \textbf{No distributed training.} Bake-off parallelism is
        limited to a single machine and $\le 4$ candidate threads.
  \item \textbf{Single-process Optuna.} The sweep runs with
        \code{n\_jobs=1} at the study level to keep bookkeeping
        simple; multi-process TPE requires shared storage.
\end{itemize}

\section{Model and Algorithm Limitations}
\label{sec:lim-models}

\begin{itemize}[leftmargin=1.4em]
  \item \textbf{Fixed candidate set.} The bake-off compares only
        RandomForest, XGBoost, LightGBM, and a linear baseline
        (LogisticRegression or Ridge). Stacking, blending, deep
        tabular models (e.g.\ TabNet, FT-Transformer), and
        CatBoost are not currently included.
  \item \textbf{No feature engineering search.} Preprocessing is a
        fixed \code{ColumnTransformer}; feature interactions,
        target encoders, and count encoders are not attempted.
  \item \textbf{No calibration or thresholding.} Classification
        outputs are point predictions; probability calibration and
        threshold tuning are not performed.
  \item \textbf{Uni-target regression / classification only.}
        Multi-label, multi-target, and hierarchical tasks are not
        supported.
\end{itemize}

\section{Agent Reliability and Hallucination}
\label{sec:lim-hallucination}

\begin{itemize}[leftmargin=1.4em]
  \item \textbf{Wrong tool selection.} The LLM may occasionally
        choose the wrong tool, e.g.\ calling
        \tool{run\_parallel\_bake\_off} without confirming the
        target column when it should have asked.
  \item \textbf{Argument fabrication.} Free-form arguments such as
        \code{target\_column} are constructed by the LLM. If the
        user's phrasing is ambiguous, the LLM may invent a plausible
        column name; this is why the modeling tool now returns a
        friendly error listing the actual columns
        (Section~\ref{sec:qual-results}).
  \item \textbf{Result misinterpretation.} The LLM may summarise
        tool results incorrectly. This is mitigated by expandable
        tool-call cards in the UI so the user can inspect raw JSON.
  \item \textbf{Iteration ceiling.} Long compound tasks may hit the
        default \code{MAX\_TOOL\_ITERATIONS=8} limit before
        completing; the platform will end with a plain-text apology
        rather than an error.
\end{itemize}

\section{Scalability Limitations}
\label{sec:lim-scale}

\begin{itemize}[leftmargin=1.4em]
  \item \textbf{Single user.} There is no user identity, no
        per-user directory, and no access control. Multi-tenant
        deployment would require reintroducing an auth layer and a
        per-user workspace.
  \item \textbf{Streamlit re-runs.} Streamlit re-runs the whole
        script on every widget event; long tool outputs are
        re-rendered on each re-run.
  \item \textbf{No horizontal scaling.} There is no queue, no
        worker process pool, and no orchestration for concurrent
        chat sessions.
\end{itemize}

\section{Security Limitations}
\label{sec:lim-security}

\begin{itemize}[leftmargin=1.4em]
  \item \textbf{No sandboxing of pandas.} Uploaded files are parsed
        with \code{pandas.read\_csv}, which is safe against arbitrary
        code execution but is exposed to CSV injection into
        downstream tools that consume string values.
  \item \textbf{Attacker-influenced tool results.} Data-derived
        strings are echoed back to the LLM and could contain
        prompt-injection payloads. Mitigation: raw tool results are
        visible in the UI; the platform is not run unattended.
  \item \textbf{Secrets in \file{.env}.} The API key is stored in
        plain text in the working directory. Users should treat
        \file{.env} as sensitive and keep it out of version
        control (the \file{.gitignore} already excludes it).
\end{itemize}

\section{Reproducibility Limitations}
\label{sec:lim-repro}

\begin{itemize}[leftmargin=1.4em]
  \item \textbf{LLM non-determinism.} Even at low temperature,
        Groq's streaming responses vary run to run; the prose is
        not byte-reproducible.
  \item \textbf{Provider drift.} A Groq-side update to
        \code{llama-3.3-70b-versatile} would change behaviour
        without any change in the local repository.
  \item \textbf{No dataset hash.} The platform does not currently
        record a hash of the uploaded dataset. Recomputing artefacts
        after a data change is silent.
\end{itemize}

\section{Dependency Footprint}
\label{sec:lim-deps}

The dependency list (Table~\ref{tab:deps}) is heavy for a research
tool: scikit-learn, XGBoost, LightGBM, SHAP, Optuna, matplotlib, and
ReportLab together account for hundreds of megabytes when installed.
This is intentional --- the platform's value proposition is a
batteries-included pipeline --- but it does mean that installation
is not instantaneous and that Docker images built naively will be
sizable.

\begin{table}[H]
\centering
\caption{Runtime dependencies of \projname{} and their purpose.
Every entry is a direct import in the code.}
\label{tab:deps}
\small
\begin{tabularx}{\linewidth}{@{}lXl@{}}
\toprule
\textbf{Package} & \textbf{Purpose in \projname} & \textbf{Category} \\
\midrule
streamlit & UI                              & UI \\
fastmcp   & MCP servers + in-process transport & Protocol \\
groq      & Chat Completions API client     & LLM \\
numpy, pandas, pyarrow & Data handling      & Data \\
scikit-learn & Preprocessing + ML pipeline  & ML \\
scipy     & Statistical helpers             & ML \\
xgboost, lightgbm & Gradient boosting       & ML \\
shap      & Explainability                  & ML \\
optuna    & Hyperparameter search           & ML \\
pandera   & Schema validation               & Data \\
matplotlib & Plot rendering                 & Viz \\
reportlab & PDF generation                  & Export \\
pydantic, pydantic-settings & Typed config  & Config \\
python-dotenv & \file{.env} loading         & Config \\
\bottomrule
\end{tabularx}
\end{table}

\section{Architectural Limitations}
\label{sec:lim-arch}

\begin{itemize}[leftmargin=1.4em]
  \item \textbf{No message persistence.} Chat history is held in
        Streamlit session state; a page refresh loses it. The
        earlier prototype (removed during migration) persisted
        conversations in a database.
  \item \textbf{No MCP resource endpoints.} The MCP protocol
        supports resources (addressable read-only content) but the
        current servers expose none.
  \item \textbf{No observability stack.} There is no Prometheus
        exporter, no tracing (OpenTelemetry), and no structured
        JSON log output. Only stderr logs are produced.
  \item \textbf{No CI/CD.} There is no automated build, lint, or
        test in the repository.
\end{itemize}

Every limitation above is addressed as future work in
Chapter~\ref{ch:conclusion}, either as a small extension or as a
clearly-labelled research direction.
